\documentclass[11pt]{article}
\usepackage[a4paper,margin=2.6cm]{geometry}
\usepackage{amsmath,amssymb,amsthm,booktabs}
\newtheorem{theorem}{Theorem}
\newtheorem{lemma}[theorem]{Lemma}
\newtheorem{proposition}[theorem]{Proposition}
\theoremstyle{remark}\newtheorem*{remark}{Remark}
\title{The sampling constant of a hole plus a sub-Nyquist prefix\\
\large One set, not $aL+bL$}
\author{D.\ Joubert}
\date{5 September 2026}
\begin{document}
\maketitle
\begin{abstract}
Under RH the floor \(c_L\) of the truncated Weil form is the lower
sampling constant of the zero set \(\Gamma\) for \(PW_{L/2}\). The
empirical law \(-\ln c_L\approx aL(\gamma_1-\nu)_++bL\sum(\mathrm{gap}-\nu)_+\)
splits that geometry into two Slepian factors and over-predicts the
narrow deserts by \(1.4\)--\(1.8\) at a common cutoff \(T_0=320\).
The honest object is one measurable set
\(E_L=(-\gamma_1,\gamma_1)\cup\bigcup_{\Delta>\nu}(\gamma_k,\gamma_{k+1})\).
Beurling still only gives \(c_L>0\). Slepian on the largest component
of \(E_L\) gives an upper bound on \(c_L\), already in
\emph{sampling-floor}. A matching lower bound that sees all of \(E_L\)
is not proved: Bernstein plus Duffin--Schaeffer loses the exponent.
Numerically \(-\ln\lambda/(L|E|)\) is constant per character and not
across characters. The missing term is therefore either a
Beurling--Malliavin density of \(\Gamma\) weighted by \(E_L\), or a
character-wise constant \(C(\chi)\) times \(|E|\). This note records
the theorems, the measurements, and where a proof has to start. It
does not move RH.
\end{abstract}

\section{What is a theorem}

Let \(Q_L\) be the truncated Weil form on even functions supported in
\([-L/2,L/2]\), and \(c_L=\inf Q_L/\|f\|^2\). Write \(\tau=L/2\),
\(\nu=2\pi/L\), \(F=\hat f\in PW_\tau\). Under RH,
\(Q_L(f)=\sum_{\gamma\in\Gamma}F(\gamma)^2\).

\begin{theorem}[floor, \emph{sampling-floor} Th.~1]
Assume RH. Then \(c_L>0\), and the finite-basis minima decrease to it.
\end{theorem}

The proof is Beurling: a separated subset of \(\Gamma\) has lower
density \(>\tau/\pi\). Density sees positivity, not the size of \(c_L\).

\begin{proposition}[desert upper bound, \emph{sampling-floor} Prop.~2]
Assume RH. Then
\(c_L\le C_\tau(1+\log\gamma_1)\,(1-\lambda_0(\tau(\gamma_1-1)))\),
hence \(-\ln c_L\ge L(\gamma_1-1)-O(\log)\).
\end{proposition}

This is Slepian on \((-\gamma_1,\gamma_1)\). At \(\mu=11\) for \(\zeta\)
it accounts for \(28\%\) of the measured depth. Replacing the desert
by the union \(E_L\) and computing \(1-\lambda_0(E_L,\tau)\) does not
bound \(c_L\) in either direction (factor \(25\) one way, \(10^{16}\)
the other, \emph{sampling-floor} §4).

\section{The set}

\[
E_L=(- \gamma_1,\gamma_1)\cup\bigcup_{\gamma_{k+1}-\gamma_k>\nu}(\gamma_k,\gamma_{k+1}).
\]
Landau's count on this set is \(n_L(E)=(\tau/\pi)|E|\). The number of
endpoints of the components of \(E\cap\mathbb R_+\) is
\(n_\partial=1+2N_{\mathrm{long}}\). Interlacing of
\(A=\chi_E P_\tau\chi_E\) restricted to
\(V=\{F\in PW_\tau:F|_{\partial E}=0\}\) gives
\[
\dim_\delta(E)=\#\{\lambda_j(A|_V)\ge 1-\delta\}=n_L(E)-n_\partial+O(\log).
\]
That count is proved (\texttt{report/one-set-lemmas.md}). It is not
an exponent.

\section{What the hold-out did}

Sixteen primitive \(L\)-functions have complete one-sided zero lists
to \(T_0=320\). Coefficients \(a=1.71\), \(b=0.97\) are those of
\(\zeta\) alone. Three characters measured after the definition:

\begin{center}\small
\begin{tabular}{lccc}\toprule
\(\chi\) & \(s_{\mathrm{hat}}\) & \(s_{\mathrm{pred}}(a,b)\) & ratio\\\midrule
\(\chi_{29}\), \(11\)--\(22\) & \(0.390\) & \(0.555\) & \(0.70\)\\
\(\chi_{17}\), \(22\)--\(74\) & \(0.728\) & \(1.187\) & \(0.61\)\\
\(\chi_5\), \(30\)--\(50\) & \(2.019\) & \(3.048\) & \(0.66\)\\\bottomrule
\end{tabular}
\end{center}

Narrow deserts are over-predicted by \(1.4\)--\(1.8\). The split
\(aL+bL\) charges every long gap as a second independent Slepian in
the same measure \(E_L\).

\section{Three scales against the measurements}

On the same lists, with \(\ell=-\ln\lambda_{\min}\):

\begin{center}\small
\begin{tabular}{lcccc}\toprule
\(\chi\) & \(\mu\) & \(\ell/(L|E|)\) & \(\ell/(L|I_{\max}|)\) & \(\ell/\dim\)\\\midrule
\(\chi_5\) & \(38\)--\(62\) & \(0.198\)--\(0.199\) & \(3.6\)--\(4.6\) & \(2.4\)--\(2.8\)\\
\(\chi_{17}\) & \(11\)--\(74\) & \(0.11\)--\(0.21\) & \(0.55\)--\(3.1\) & \(1.66\)--\(2.0\)\\
\(\chi_{29}\) & \(11\)--\(74\) & \(0.082\)--\(0.095\) & \(0.14\)--\(1.75\) & \(1.13\)--\(1.59\)\\\bottomrule
\end{tabular}
\end{center}

\(L|I_{\max}|\) is not the scale of \(\ell\). \(L|E|\) is constant
\emph{inside} a character and not across characters (factor \(2.5\)).
A universal lower bound \(\ell\ge 0.08\,L|E|\) survives the table and
has the same slack as \(a+b\). \(\dim\) is the least bad single
predictor and is not absolute. Details:
\texttt{report/one-set-ratios.md}.

\section{What a lower bound would have to be}

The inequality that would replace \(aL+bL\) is: for every
\(F\in PW_\tau\),
\[
\sum_{\gamma\in\Gamma}|F(\gamma)|^2
\;\ge\;
\exp\bigl(-C\,\tau\,|E_L|\bigr)\,\|F\|^2
\]
with \(C\) absolute. Slepian gives the integral of \(|F|^2\) off
\(I_{\max}\), not the sum over \(\Gamma\), and not off all of \(E_L\).
Passing from the integral to the sum by Bernstein
(\(\|F'\|\le\tau\|F\|\)) and Duffin--Schaeffer on each short gap makes
the bracket negative of order \(e^{-L\gamma_1}\)
(\emph{sampling-floor} §4). That route is closed.

A form \(c_L\ge\exp(-C\dim(E))\) is false as soon as
\(\dim\ll L|I_{\max}|\) (\(\chi_{29}\) at \(\mu=11\): \(1.1\) against
\(4\)). A form \(c_L\ge\exp(-L|I_{\max}|)\) is false as a lower bound
(\(\zeta\) at \(\mu=11\): \(e^{-34}\) against measured \(e^{-110}\));
it is the desert \emph{upper} bound on \(c_L\), already known.

\section{Where a proof starts}

Two objects, neither of them a scan of \(Q\).

\paragraph{Weighted Beurling--Malliavin.}
The multiplier radius of \(\Gamma\) for type \(\tau\), computed with
a hole kernel \(1_{E_L}\). If that radius is \(C|E_L|\), the
inequality of §5 is a theorem. This is an article in harmonic
analysis.

\paragraph{A profile, not a length.}
The bottom eigenvector \(v_0\) of \(Q\) is known numerically. The
mass \(|\hat v_0(\gamma)|^2\) and the derivatives
\(\partial(-\ln\lambda)/\partial\gamma_k\) on the zero-side Gram say
whether the cost is one gap, an energy of gaps, or a factor
\(C(\chi)\) independent of \(L\). The ratios \(\ell/(L|E|)\) already
point to \(C(\chi)\). That computation is the sequel to this note.

\section*{Status}
Theorem 1 and Proposition 2 are unchanged. The set \(E_L\) and the
count \(\dim(E)\) are defined and measured. The law \(aL+bL\) is
diagnosed, not refit. No lower bound that sees \(E_L\) as a set is
proved. Nothing bears on RH.
\end{document}
